\documentclass[conference]{IEEEtran}
\usepackage{cite}
\usepackage{graphicx}
\usepackage{amsmath,amssymb}
\usepackage{url}
\usepackage{booktabs}

\begin{document}

\title{Dual-UPF Traffic Steering for 5G Network Slices in Stadium and Mass Event Deployments}

\author{
\IEEEauthorblockN{Ajay Chander Ravichandran}
\IEEEauthorblockA{\textit{IMT Atlantique}\\Rennes, France\\ajay.ravichandran@imt-atlantique.net}
\and
\IEEEauthorblockN{Bruno Stevant}
\IEEEauthorblockA{\textit{IMT Atlantique}\\Rennes, France\\bruno.stevant@imt-atlantique.fr}
\and
\IEEEauthorblockN{Hadi Houssainy}
\IEEEauthorblockA{\textit{IMT Atlantique}\\Rennes, France\\hadi.houssainy@imt-atlantique.fr}
}

\maketitle

\begin{abstract}
This paper addresses the challenge of optimal User Plane Function (UPF) selection in high-density 5G stadium environments characterized by extreme traffic spikes and heterogeneous service requirements. We propose a slice-aware traffic steering framework based on Evolutionary Game Theory (EGT) that dynamically balances load between Edge (MEC) and Central Cloud (CC) UPFs. The system is implemented and validated on an OpenAirInterface (OAI) 5G Core emulation testbed, scaled to 1,000 concurrent physical UEs with differentiated eMBB (80\%) and URLLC (20\%) profiles. By leveraging a high-performance Go-based user-plane multiplexing engine, we achieve 100\% PDU session establishment success at scale. Our results demonstrate that the EGT-based steering ensures 100\% QoS compliance for eMBB traffic (300 ms PDB) during 5.5x peak halftime surges, while providing a fair distribution that maximizes resource utilization. We further analyze the infrastructure-limited boundaries for strict URLLC constraints (10 ms PDB), providing a blueprint for reliable edge service delivery in mass event deployments.
\end{abstract}

\begin{IEEEkeywords}
5G network slicing, User Plane Function (UPF) selection,
multi-access edge computing (MEC), quality of service (QoS),
radio access network (RAN), evolutionary game theory (EGT),
OpenAirInterface
\end{IEEEkeywords}

% -------------------------------------------------------
\section*{Abbreviations}
\begin{center}
\renewcommand{\arraystretch}{1.1}
\begin{tabular}{ll}
\hline
\textbf{Abbreviation} & \textbf{Full Form} \\
\hline
5G    & Fifth Generation \\
eMBB  & Enhanced Mobile Broadband \\
EGT   & Evolutionary Game Theory \\
MEC   & Multi-Access Edge Computing \\
OAI   & OpenAirInterface \\
PDB   & Packet Delay Budget \\
QoS   & Quality of Service \\
RAN   & Radio Access Network \\
SMF   & Session Management Function \\
UE    & User Equipment \\
UPF   & User Plane Function \\
URLLC & Ultra-Reliable Low-Latency Communication \\
CUPS  & Control and User Plane Separation \\
SBA   & Service-Based Architecture \\
BLSTM & Bidirectional Long Short-Term Memory \\
PFCP  & Packet Forwarding Control Protocol \\
PDU   & Protocol Data Unit \\
NR    & New Radio \\
NF    & Network Function \\
ILP   & Integer Linear Programming \\
OPEX  & Operational Expenditure \\
GTP   & GPRS Tunnelling Protocol \\
NSSF  & Network Slice Selection Function \\
AMF   & Access and Mobility Management Function \\
NRF   & Network Repository Function \\
UDM   & Unified Data Management \\
AUSF  & Authentication Server Function \\
UDR   & Unified Data Repository \\
\hline
\end{tabular}
\end{center}
% -------------------------------------------------------

\section{Introduction}

Stadiums and large-scale public events are increasingly becoming
testbeds for next-generation connectivity. As 5G NR deployments
mature, the distinct service classes they support open up
possibilities that were simply not feasible with previous network
generations: eMBB for high-throughput fan engagement, and URLLC for
mission-critical operations such as drone coordination. 
Delivering immersive AR/VR experiences to
tens of thousands of spectators simultaneously, while also
supporting real-time safety systems and pervasive crowd monitoring,
places extraordinary demands on network infrastructure. Latencies
must remain below 20~ms for time-sensitive applications, throughput
must scale to gigabit levels, and reliability cannot drop below
99.999\% for safety-critical slices. Meeting all of these
requirements in a single venue, at peak load, is a genuinely hard
problem.

A key part of the solution lies at the edge. By deploying MEC
nodes close to the RAN, operators can offload compute-intensive
tasks to local servers, including 360\textdegree{} video
processing, AI-based crowd analytics, and real-time resource
orchestration. Moving this processing away from the core cuts
round-trip latency significantly and improves the Quality of
Experience (QoE) for end users. The growing availability of
localised UPFs, driven in part by the broader expansion of
distributed data centres, makes this architecture increasingly
practical. However, relying on a single centralised UPF in a
high-density scenario quickly becomes a bottleneck: bandwidth
gets congested, propagation delays creep up, and the QoS
guarantees associated with each network slice begin to erode.
A multi-UPF deployment addresses this by steering traffic to
whichever edge node is best placed to serve a given UE,
provided the selection logic is sound. Getting that logic right
requires algorithms that can jointly account for load, latency,
per-slice QoS requirements, and fairness across UEs, so that
neither eMBB fans nor URLLC safety services are starved of
resources. This paper proposes and evaluates such an algorithm,
with the aim of enabling reliable 5G edge service delivery across
all slice types in mass event environments.

\section{Related Work}

The transition from hardware-centric legacy infrastructure to the
5G System (5GS) is fundamentally shaped by the adoption of a
Service-Based Architecture (SBA). A cornerstone of this shift is
Control and User Plane Separation (CUPS), which decouples session
management logic from data-plane routing, enabling the UPF to be
instantiated anywhere in the topology---from centralised data
centres to the network edge. This architectural flexibility
transforms UPF selection from a static assignment into a dynamic,
slice-aware optimisation problem, and provides the architectural
motivation for the work surveyed in this section.

The body of work relevant to this paper spans three interconnected
areas: 5G network slicing and resource management, UPF placement
and selection in mobile core networks, and MEC deployment in
high-density environments.

Network slicing has been extensively studied since its introduction
as a core concept in 3GPP Release~15. Early work focused on the
radio access side, with proposals for slice-aware scheduling and
resource partitioning at the base station
level~\cite{zhang2019network}. More recent contributions have
addressed end-to-end slice management, including control-plane
signalling, isolation guarantees, and dynamic admission control
across eMBB and URLLC traffic
classes~\cite{popovski2018wireless}. These slices impose distinct
KPI targets: eMBB demands peak data rates of 20\,Gbps downlink
and 10\,Gbps uplink; URLLC mandates end-to-end latency below
1\,ms with $99.999\%$ availability~\cite{3gpp_ts22261}.
While these works establish the slicing framework, they largely
assume a fixed core network topology and do not consider UPF
placement as a variable in the optimisation. More recently,
distributed QoS-aware slicing algorithms have been extended beyond
5G; \cite{Asheralieva2024slicing} proposes an ADMM-based framework
combining graph neural networks with distributed optimisation to
achieve provably convergent resource slicing under dynamic and
unreliable network conditions.

UPF placement and traffic steering have received growing attention
as operators move toward distributed core architectures. Several
studies have framed UPF placement as a facility location problem,
optimising for latency or load under capacity
constraints~\cite{Goshi2024upf}. Others have proposed dynamic UPF
migration triggered by mobility patterns or traffic
shifts~\cite{sama2016software}. However, most of these approaches
target general LTE/5G deployments and do not account for the
extreme density and heterogeneous QoS profiles characteristic of
mass event scenarios, where eMBB and URLLC traffic coexist
and compete for the same edge infrastructure. A promising
EGT-based approach for dynamic UPF selection, where UEs adapt
their choice to minimise end-to-end delay using replicator dynamics
converging to a stable equilibrium, is proposed in
\cite{Alevizaki2021upf}. Complementing this, \cite{LeyvaUPC2023}
proposes an ML-based framework that proactively schedules UPF
placement and chaining reconfigurations by predicting latency
violations, reducing unnecessary reconfiguration events while
keeping QoS within bounds the majority of the time.

MEC deployment in high-density venues has been explored in the
context of sports stadiums and concerts, with work addressing task
offloading, video caching, and AR/VR
streaming~\cite{mehrabi2019device}. Notably, Chochliouros
et~al.\ validate a two-tier MEC architecture with
Cloud-Enabled Small Cells at a live football stadium,
demonstrating that local edge processing effectively reduces
backhaul load in high-density
venues~\cite{Chochliouros2020stadium, Kostopoulos2019edge}. These
studies demonstrate the latency and throughput benefits of edge
processing but typically treat the UPF as a fixed anchor point
rather than a selectable network element. The interaction between
MEC service placement and UPF selection under slice-specific QoS
constraints remains largely unaddressed in the literature.

This paper bridges these gaps by jointly considering UPF selection
and slice-aware resource allocation in a dual-UPF, single-MEC
architecture, evaluated specifically in the context of mass event
deployments where traffic heterogeneity and peak-load conditions
make the selection problem non-trivial.

\section{Methodology}
\label{sec:methodology}

%--------------------------------------------------------------
\subsection{System Model and Network Topology}
%--------------------------------------------------------------

The network is modelled as a directed graph $\mathcal{G}
(\mathcal{N},\mathcal{E})$, where $\mathcal{N}$ denotes
physical nodes and $\mathcal{E}$ directed
links~\cite{LeyvaUPC2023}. The node set comprises
stadium-side gNBs forming the 5G NR RAN, a MEC edge server
$c_\text{mec} \in \mathcal{N}_c$ at a stadium aggregation
point, and a central cloud server $c_\text{cc} \in
\mathcal{N}_c$ in the operator core. Each server $c$ has
maximum processing capacity $C_c$. Each link $(u,v) \in
\mathcal{E}$ carries a bandwidth capacity $B_{u,v}$
and one-way delay $d_{u,v}$ comprising propagation and
switching components.
$\text{UPF}_\text{mec}$ provides low propagation latency
at constrained processing capacity.
$\text{UPF}_\text{cc}$ provides higher processing headroom
at greater propagation cost. Propagation delay is modelled
at $5~\mu\text{s/km}$ of optical fibre~\cite{Alevizaki2021upf},
and the central UPF is sited such that
$t_{\text{prop}_\text{cc}} / t_{\text{prop}_\text{mec}} = 4$.

Per 3GPP~TS~23.501~\cite{3gpp_ts23501}, UPF selection and
session management are the responsibility of the SMF.
In this work the EGT algorithm of
Section~\ref{sec:egt} runs within the SMF as a
population-aware steering policy, maintaining population
state vectors $\mathbf{x}^g(t)$ and issuing anchor
assignments to UEs via PFCP over the N4 reference
point~\cite{3gpp_ts29244}.

\begin{figure}[htbp]
    \centering
    \includegraphics[width=0.45\textwidth]{image1.png}
    \caption{5G multi-UPF network architecture. eMBB and URLLC
    UEs connect to the gNB via the Uu air interface. The gNB
    forwards traffic over N3 to either the MEC edge UPF or the
    central cloud UPF. The SMF steers sessions between both UPFs
    via the N4 control interface.}
    \label{fig:network_architecture}
\end{figure}

As illustrated in Fig.~\ref{fig:network_architecture}, the
proposed topology differentiates between eMBB and URLLC
sessions. Traffic originating from UEs is carried via the
Uu interface to the gNB, then steered over the N3 user-plane
interface to either $\text{UPF}_\text{mec}$ or
$\text{UPF}_\text{cc}$. Session management and UPF selection
are governed by the SMF via the N4/PFCP interface.

The UE population is partitioned into two slice groups
$g \in \{1,2\}$. We explicitly exclude mMTC (massive Machine Type 
Communication) traffic from the evaluation, as its low-throughput, 
highly bursty nature exerts negligible selection pressure on UPF 
processing capacity compared to broadband video and spatial rendering.
Group $g=1$ comprises eMBB users
(S-NSSAI: SST=1) who represent the dominant traffic
class at a stadium. eMBB traffic in this scenario
primarily consists of HD video replay consumption,
social media upload, and venue app usage. Active
HD video streaming demands approximately $25$~Mbps per
device~\cite{yaqoob2020survey_360video}, while the
average per-device load in a dense shared-cell
environment, as observed in large-event
deployments, is significantly lower due to scheduler
sharing across concurrent
users~\cite{superbowl_lvii,ericsson_f1_2023}.
We model $\rho_\text{eMBB} = 10$~Mbps as the mean
offered load per eMBB UE under peak conditions, which
is conservative relative to peak single-device rates
but reflects realistic concurrent cell loading.

Group $g=2$ comprises URLLC users (S-NSSAI: SST=2)
supporting latency-sensitive AR/VR overlay services
offloaded to the MEC node. In this deployment, URLLC
sessions serve real-time AR rendering tasks in which
compute-intensive frame processing is handled at the
MEC server and only compressed output is delivered
to headsets, requiring $\rho_\text{URLLC} = 25$~Mbps
per session~\cite{yaqoob2020survey_360video} with a
strict packet delay budget. Both the per-slice
packet delay budgets $\text{PDB}_g$ and 5QI values
are derived from Table~5.7.4-1 of
3GPP~TS~23.501~\cite{3gpp_ts23501} and
TS~22.261~\cite{3gpp_ts22261}: $\text{PDB}_1 = 300$~ms
for eMBB video (5QI~6, non-GBR) and
$\text{PDB}_2 = 10$~ms for interactive AR/VR
(5QI~83, Delay-Critical GBR)~\cite{3gpp_ts22261}.
All parameters are listed in Table~\ref{tab:params}.

\begin{table}[ht]
\centering
\caption{System, slice, and algorithm parameters}
\label{tab:params}
\renewcommand{\arraystretch}{1.15}
\begin{tabular}{p{1.8cm}p{1.9cm}p{3.9cm}}
\toprule
\textbf{Symbol} & \textbf{Value} &
\textbf{Description and source} \\
\midrule
\multicolumn{3}{l}{\textit{Network topology}} \\
$C_c^\text{mec}$ & 40 cores (2.4 GHz) &
  MEC server capacity (edge-constrained) \\
$C_c^\text{cc}$ & 120 cores (2.4 GHz) &
  Central server capacity \\
$B_{u,v}$ & 100 Gbps &
  Link bandwidth \\
$t_{\text{prop}_\text{mec}}$ & 0.25 ms &
  Propagation delay to local UPF \\
$t_{\text{prop}_\text{cc}}$ & 1.0 ms &
  Propagation delay to central UPF \\
\midrule
\multicolumn{3}{l}{\textit{Slice parameters
  (3GPP TS 22.261~\cite{3gpp_ts22261},
   TS 23.501~\cite{3gpp_ts23501})}} \\
$M_1$ & 800 &
  eMBB UE population (SST=1) \\
$M_2$ & 200 &
  URLLC UE population (SST=2) \\
$\rho_\text{eMBB}$ & 10 Mbps &
  Mean per-UE eMBB load at
  peak~\cite{superbowl_lvii,ericsson_f1_2023} \\
$\rho_\text{URLLC}$ & 25 Mbps &
  Per-session AR/VR offload
  load~\cite{yaqoob2020survey_360video} \\
$\text{PDB}_1$ & 300 ms &
  eMBB PDB, 5QI~6~\cite{3gpp_ts23501} \\
$\text{PDB}_2$ & 10 ms &
  URLLC AR/VR PDB,
  5QI~83~\cite{3gpp_ts22261} \\
\midrule
\multicolumn{3}{l}{\textit{UPF delay
  model~\cite{Alevizaki2021upf}}} \\
$k_\text{mec}$ & $4.833\!\times\!10^{-4}$ &
  Processing exponent, local UPF
  (units: Mbps$^{-1}$); calibrated via
  grid search against Fig.~2
  of~\cite{Alevizaki2021upf} \\
$k_\text{cc}$ & $4.833\!\times\!10^{-5}$ &
  Processing exponent, central UPF;
  ratio $k_\text{mec}/k_\text{cc}=10$
  per~\cite{Alevizaki2021upf} \\
\midrule
\multicolumn{3}{l}{\textit{EGT algorithm~\cite{Alevizaki2021upf}}}\\
$\beta$ & 1 &
  Payoff scaling constant \\
$a$ & 1 &
  Review rate upper bound \\
$\varepsilon$ & 0.01 &
  Payoff convergence threshold \\
$\Delta t$ & 0.05 &
  Discrete replicator step size \\
\bottomrule
\end{tabular}
\end{table}

%--------------------------------------------------------------
\subsection{Traffic Dynamics at Mass Events}
\label{sec:traffic}
%--------------------------------------------------------------

A defining characteristic of the stadium scenario is the
pronounced non-stationarity of traffic demand. Real-world
measurements at large events document this pattern
clearly. At Super Bowl LIII, 24~TB of data were
transferred within the stadium alone, with the data
network spiking to 13.06~Gbps during the halftime show
from over 30,000 concurrent mobile
devices~\cite{nokia_live_events}. At Super Bowl~LVII,
a peak Wi-Fi throughput of 25.8~Gbps was recorded,
and the halftime show drove a 450\% increase in
bandwidth usage relative to the pre-show
baseline~\cite{superbowl_lvii}. At the 2023 Formula~1
Grand Prix in Mexico~City, 5G accounted for 40\% of
over 117~TB of total mobile data traffic, with average
5G download speeds of 274~Mbps for connected
users~\cite{ericsson_f1_2023}. These measurements
establish three distinct load phases that our model
must capture. The halftime demand spike in our
evaluation is modelled at $\lambda = 5.5\times$ the
pre-spike baseline, where a 450\% increase above the 
pre-match baseline corresponds to a load multiplier 
of $\lambda = 1 + 4.5 = 5.5$.
At this multiplier the total offered load is
$(M_1\rho_\text{eMBB} + M_2\rho_\text{URLLC})\,\lambda
= 13~\text{Gbps} \times 5.5 = 71.5~\text{Gbps}$.

\textbf{Pre-match ramp-up (Phase~I).} Spectators enter
the venue over 30--60 minutes. eMBB sessions increase
monotonically as users connect, stream highlights, and
access venue applications. Load is moderate and
$\text{UPF}_\text{mec}$ can absorb the majority of
sessions with favourable delay; the EGT equilibrium
concentrates sessions at the edge node.

\textbf{Peak-load phase (Phase~II).} At kick-off and
halftime, concurrent eMBB and URLLC sessions generate
the aggregate demand surge documented in field
measurements. Phase II spans 45 minutes, evaluated at 
3 simulation steps per minute, yielding a 135-step 
evaluation window. The exponential processing delay model in
Eq.~\eqref{eq:upf_delay_mec}--\eqref{eq:upf_delay_cc}
captures the rapid onset of queuing congestion at
$\text{UPF}_\text{mec}$ as aggregate load approaches
capacity. Rising processing delay reduces the payoff of
the local UPF below the population average, and the
replicator dynamics redistribute sessions to
$\text{UPF}_\text{cc}$. This phase creates the
conditions under which the EGT equilibrium may
transiently lag behind rapid load shifts.

\textbf{Post-match dispersal (Phase~III).} Session
counts fall as spectators leave. The equilibrium
reverts toward $\text{UPF}_\text{mec}$, but active
URLLC sessions must not be forcibly re-anchored
during dispersal, as unnecessary PDU session
modifications introduce signalling overhead and
momentary service interruption.

%--------------------------------------------------------------
\subsection{Justification for Evolutionary Game Theory}
\label{sec:egt_justify}
%--------------------------------------------------------------

The UPF session-steering problem in a stadium environment
has four properties that make EGT the appropriate
decision framework.

\textbf{Centralised SMF Steering.} Unlike standard EGT 
formulations where individual UEs selfishly optimise their 
choices~\cite{Alevizaki2021upf}, we implement the replicator 
dynamics centrally within the SMF to compute optimal population 
fractions. This prevents the massive control-plane signalling 
overhead that would occur if thousands of UEs independently 
triggered session renegotiations, while preserving the mathematically 
stable load-balancing properties of the Nash equilibrium.

\textbf{Population-scale decentralisation.} At peak
load, up to $M_1 + M_2 = 1{,}000$ PDU sessions must be
steered concurrently. Centralised formulations such as
integer linear programming (ILP) are computationally
intractable at this scale and under dynamic
demand~\cite{LeyvaUPC2023}. EGT operates on aggregate
population fractions $\mathbf{x}^g(t)$ rather than
individual sessions, reducing the SMF's per-iteration
complexity from $O(M_g)$ to $O(|\mathcal{S}_g|) = O(2)$
for a two-UPF strategy space~\cite{Alevizaki2021upf}.

\textbf{Non-stationarity.} The three-phase load profile
changes continuously. Unlike static assignment policies,
which require offline re-optimisation to respond to load
changes, the replicator dynamics adapt $\mathbf{x}^g(t)$
continuously to the prevailing payoff landscape without
requiring a global re-solve.

\textbf{Convergence to a stable equilibrium.} The
replicator dynamics converge to a Nash equilibrium at
which all active strategies within each group achieve
equal payoff~\cite{Alevizaki2021upf}. At equilibrium,
neither UPF is overloaded relative to the other given
the current population mix. The current payoff function
(Eq.~\eqref{eq:payoff}) optimises end-to-end delay
rather than PDB margin directly; the relationship between
equilibrium delay and QoS compliance is discussed in
Section~\ref{sec:results}.

\textbf{Low N4 signalling overhead.} The SMF maintains
population state internally; no per-session global state
exchange is required. The algorithm issues N4/PFCP
session modification rules only for newly established
PDU sessions, consistent with
the PFCP interaction model of
3GPP~TS~29.244~\cite{3gpp_ts29244}.

%--------------------------------------------------------------
\subsection{UPF Strategy Space and Population State}
\label{sec:strategy}
%--------------------------------------------------------------

For each group $g \in \{1,2\}$, the strategy set is
$\mathcal{S}_g = \{\text{UPF}_\text{mec},
\text{UPF}_\text{cc}\}$.
The SMF maintains the population state of group $g$ as:
\begin{equation}
    \mathbf{x}^g(t) = \bigl[x_\text{mec}^g(t),\;
    x_\text{cc}^g(t)\bigr], \quad
    x_\text{mec}^g + x_\text{cc}^g = 1
    \label{eq:popstate}
\end{equation}
where $x_i^g(t)$ is the fraction of sessions in group
$g$ currently anchored at $\text{UPF}_i$, computed from
per-UPF session counts reported over N4/PFCP at each
iteration interval $\Delta t$.

%--------------------------------------------------------------
\subsection{Delay Model and Payoff Function}
\label{sec:delay}
%--------------------------------------------------------------

The end-to-end delay for a session in group $g$ anchored
at $\text{UPF}_i$ under load multiplier $\lambda$ is:
\begin{equation}
    d_i^g(\mathbf{x},\lambda) = t_{\text{prop}_i}
    + t_{\text{UPF}_i}^g(\mathbf{x},\lambda)
    \label{eq:delay}
\end{equation}
where $t_{\text{prop}_i}$ is the fixed propagation delay
to $\text{UPF}_i$ and $t_{\text{UPF}_i}^g$ is the
processing and queuing delay, modelled as an exponential
function of offered load following~\cite{Alevizaki2021upf}.

The two UPFs have different resource-sharing assumptions
that match the deployed architecture. The edge
$\text{UPF}_\text{mec}$ operates with dedicated
per-slice processing: only the sessions of group $g$
itself contribute to that group's delay at the edge.
The central $\text{UPF}_\text{cc}$ is a shared node
whose processing load is the aggregate across all groups.
While Alevizaki et al. use a single shared delay formula 
for both UPFs, we adopt a dedicated-MEC model for the edge 
UPF, in which only the sessions of group $g$ contribute to 
that group's MEC processing delay. This reflects the 3GPP 
TS 23.501 network slicing architecture, where per-S-NSSAI 
resource isolation at the edge is a core requirement for 
stadium deployments with strict URLLC service guarantees 
(Section 5.7, QoS framework)~\cite{3gpp_ts23501}. The 
central UPF remains shared across all groups, consistent 
with the reference model. Accordingly:
\begin{align}
    t_{\text{UPF}_\text{mec}}^g(\mathbf{x},\lambda)
      &= \exp\!\bigl(k_\text{mec}\cdot
         \rho_g\,M_g\,x_\text{mec}^g\,\lambda\bigr)
    \label{eq:upf_delay_mec} \\[4pt]
    t_{\text{UPF}_\text{cc}}(\mathbf{x},\lambda)
      &= \exp\!\Bigl(k_\text{cc}\cdot\lambda
         \sum_{g'=1}^{G}\rho_{g'} M_{g'}
         x_\text{cc}^{g'}\Bigr)
    \label{eq:upf_delay_cc}
\end{align}
where $\rho_g$ is the per-UE offered load for slice group
$g$, $M_g$ is the group population, and $k_i$ is a
node-specific exponent whose units are Mbps$^{-1}$.
Note that $t_{\text{UPF}_\text{cc}}$ carries no group
superscript because the central UPF delay is identical
for all groups at a given load point.

The exponential form captures the superlinear onset of
queuing delay as utilisation approaches the node's
processing capacity. Since $k_\text{mec} \gg k_\text{cc}$,
$\text{UPF}_\text{mec}$ degrades faster under load than
$\text{UPF}_\text{cc}$, reproducing the edge--core
capacity asymmetry characteristic of stadium
deployments~\cite{nokia_live_events}.

The node-specific exponents are calibrated via grid search
to reproduce the reference equilibrium fractions of Fig.~2
in~\cite{Alevizaki2021upf}, using the reference parameters
$M_1=130$, $M_2=70$, $\rho=100$~Mbps. The calibrated
values $k_\text{mec} = 4.833\times10^{-4}$~Mbps$^{-1}$
and $k_\text{cc} = 4.833\times10^{-5}$~Mbps$^{-1}$
reproduce the published result to within 0.1\%:
$G_1\text{@MEC} = 17.82\%$ (reference: $\approx 17.8\%$)
and $G_2\text{@MEC} = 31.97\%$ (reference: $\approx 32.0\%$)
in 231 iterations. The 10:1 exponent ratio is preserved
per~\cite{Alevizaki2021upf}. Empirical validation against
physical testbed measurements is left for future work.
We note that these $k$ parameters, originally derived from 
optical switching hardware, are utilized here as generic capacity 
constraints to preserve the 10:1 cloud-to-edge ratio and evaluate 
the steering logic's stability. In a physical software UPF deployment, 
extreme overload would likely result in packet drops and session 
rejections rather than unbounded exponential queuing delay.

The payoff for sessions in group $g$ anchored at
$\text{UPF}_i$ is the inverse end-to-end delay:
\begin{equation}
    u_i^g(\mathbf{x},\lambda) =
    \frac{1}{d_i^g(\mathbf{x},\lambda)}
    \label{eq:payoff}
\end{equation}
It is important to note that the payoff function $u_i^g = 1/d_i^g$ is slice-agnostic; it mathematically minimises delay but is fundamentally blind to the discrete PDB thresholds. Consequently, the controller cannot prioritise URLLC over eMBB if doing so would destabilise the Nash equilibrium.

The instantaneous violation count $L_t$ is reported as
a QoS compliance metric and is intended to serve as an
input to the ML trigger described in
Section~\ref{sec:future}, but does not feed back into
the payoff function or alter the steering logic in the
current implementation.

%--------------------------------------------------------------
\subsection{EGT Replicator Dynamics}
\label{sec:egt}
%--------------------------------------------------------------

The SMF adapts session-to-UPF assignments by running the
EGT replicator dynamics of~\cite{Alevizaki2021upf}.
The review rate for sessions in group $g$ at
$\text{UPF}_i$ is:
\begin{equation}
    r_i^g(\mathbf{x},\lambda) =
    a - \beta\,u_j^g(\mathbf{x},\lambda),
    \quad \beta > 0,\quad
    \frac{a}{\beta} > u_i^g(\mathbf{x},\lambda)
    \label{eq:revrate}
\end{equation}
where $u_j^g(\mathbf{x},\lambda)$ denotes the payoff of an alternative UPF strategy $j \in \mathcal{S}_g$ chosen at random by the reviewing agent. Sessions with high payoff (low delay) are reviewed less
frequently, stabilising their assignment. Population
fractions evolve according to:
\begin{equation}
    \dot{x}_i^g(t) = \beta\Bigl(u_i^g(\mathbf{x},\lambda)
    - \bar{u}^g(\mathbf{x},\lambda)\Bigr)x_i^g, \quad
    \bar{u}^g = \sum_i x_i^g\,u_i^g(\mathbf{x},\lambda)
    \label{eq:replicator}
\end{equation}
The continuous-time system converges to the interior
Nash equilibrium $u_\text{mec}^g = u_\text{cc}^g$ for
each group $g$~\cite{Alevizaki2021upf}.

In the discrete-time implementation, the SMF integrates
Eq.~\eqref{eq:replicator} with step size $\Delta t = 0.05$
up to a maximum of 2{,}000 iterations per convergence
call. Two stopping criteria are applied in conjunction:
(i) the maximum per-group payoff deviation falls below
$\varepsilon = 0.01$, i.e.\
$\max_i |u_i^g - \bar{u}^g| < \varepsilon$ for all
interior-strategy groups; and (ii) the maximum
per-element strategy change $\max_{g,i}|x_i^g{}^{(t+1)}
- x_i^g{}^{(t)}| < 2\times10^{-5}$, which detects
convergence to a boundary equilibrium where one UPF
carries zero sessions. Groups with all sessions on a
single UPF (a boundary state) are skipped in the payoff
check, as the payoff equality condition is not applicable
at the boundary.

Established sessions are not re-anchored mid-flight
within an iteration to avoid service interruption; the
redistribution applies to new establishments and
re-establishments only. 

\section{System Implementation and Architectural Considerations}
\label{sec:implementation}

To validate the EGT steering logic under high load, we
developed an emulation testbed using OpenAirInterface
(OAI) and a modified version of the \texttt{gnbsim}
emulator. Supporting 1,000 concurrent PDU sessions
required architectural changes to the standard
per-UE processing model.

\subsection{High-Performance User-Plane Multiplexing}

Standard mobile simulators typically spawn independent
goroutines for each User Equipment, leading to resource
contention at high densities. We addressed this by
re-engineering the emulator's user-plane architecture
into a multiplexed engine using shared
\texttt{sharedEncap} and \texttt{sharedDecap} loops with
$O(1)$ lookup maps---specifically \textit{IpToCamper}
and \textit{TeidToCamper}---to route traffic. This
allows the system to identify the correct UE context
and GTP tunnel parameters for any incoming packet in
constant time, enabling the testbed to sustain the
session volume required by the stadium scenario.
It should be noted that all 1,000 UE sessions are
multiplexed through a single virtual gNB instance;
this differs from a physical multi-cell deployment and
constitutes a simplification in the emulation
fidelity.

\subsection{EGT Controller Integration}

The EGT controller retrieves session counts via the N4/PFCP 
interface using Session Report messages and utilizes the O1 
interface for management-plane telemetry. At each iteration, 
the controller attempts to retrieve live session counts from 
the SMF to update the population state $\mathbf{x}^g(t)$. 
When live SMF data is available, this ensures the population 
state reflects actual PDU session anchor distribution rather 
than purely mathematical projection. In scenarios where the
SMF interface was not reachable during evaluation, the
controller operated in model-only mode, and results
are reported accordingly.

\section{Results and Discussion}
\label{sec:results}

The EGT steering controller is evaluated across four traffic
scenarios (S1--S4) against a static 50/50 baseline partition.
All scenarios use the dedicated-UPF delay model and pure
replicator dynamics of Alevizaki~et~al.~\cite{Alevizaki2021upf}. The emulation successfully provisioned 1,000 physical UE contexts within the OAI core with a 100\% PDU session establishment rate.
At peak Phase~II load ($5.5\times$ multiplier), the total offered
load is $(800 \times 10 + 200 \times 25) \times 5.5 = 71.5$~Gbps.

%--------------------------------------------------------------
\subsection{Model Validation Against Reference Equilibrium}
%--------------------------------------------------------------

Before evaluating stadium scenarios, the replicator dynamics
are validated against Fig.~2 of Alevizaki~et~al.~\cite{Alevizaki2021upf}
using the reference parameters ($M_1=130$, $M_2=70$,
$\rho=100$~Mbps). The simulator converges in 231~iterations to:
\begin{equation*}
  G_1\text{@MEC} = 17.82\%,\quad G_2\text{@MEC} = 31.97\%
\end{equation*}
against reference targets of $\approx 17.8\%$ and $\approx 32.0\%$
respectively, confirming that the calibrated parameters
$k_\text{mec} = 4.833\times10^{-4}$ and $k_\text{cc} = 4.833\times10^{-5}$
faithfully reproduce the published result.

%--------------------------------------------------------------
\subsection{Multi-Scenario Performance at Peak Load}
%--------------------------------------------------------------

Table~\ref{tab:stadium_results} summarises the key metrics extracted
at the peak load point of each scenario. The QoS compliance columns
are reported separately for eMBB (PDB~=~300~ms) and URLLC
(PDB~=~10~ms). The ``Transient Iters'' column reports the number
of replicator iterations required to reach the stopping criterion
from the scenario's initial state, providing a measure of
convergence speed and recovery robustness.

\begin{table*}[t]
\centering
\caption{EGT Steering Performance at Peak Phase~II Load (1,000 UEs).
         MEC and CC loads include the load multiplier. eMBB and URLLC
         violations are cumulative UE-steps over the 135-step Phase~II window.
         Transient~Iters for S1, S3, S4 is the static cold-start convergence
         count; for S2 it is the per-step tracking iteration count
         (warm-start dynamic mode, see Section~\ref{sec:key_findings}).}
\label{tab:stadium_results}
\renewcommand{\arraystretch}{1.15}
\begin{tabular}{lccccccc}
\toprule
\textbf{Scenario} &
  \textbf{$x_\text{mec}^\text{eMBB}$ (eq.)} &
  \textbf{$x_\text{mec}^\text{URLLC}$ (eq.)} &
  \textbf{MEC Load} &
  \textbf{CC Load} &
  \textbf{eMBB Vio.} &
  \textbf{URLLC Vio.} &
  \textbf{Transient Iters} \\
\midrule
S1 (Standard, $5.5\times$)
  & 0.156 & 0.222 & 12.9~Gbps & 58.6~Gbps & 0 & 26{,}316 & 505 \\
S2 (Extreme, $10\times$)$^\dagger$
  & 0.147 & 0.240 & 23.8~Gbps & 106.2~Gbps & 13{,}491 & 27{,}064 & 147 \\
S3 (MEC-Biased Initialisation, $5.5\times$)
  & 0.156 & 0.223 & 12.9~Gbps & 58.6~Gbps & 0 & 26{,}514 & 670 \\
S4 (CC-Biased Initialisation, $5.5\times$)
  & 0.128 & 0.198 & 11.2~Gbps & 60.3~Gbps & 0 & 44{,}027 & 270 \\
\midrule
Static 50/50 Baseline
  & 0.500 & 0.500 & 35.8~Gbps & 35.8~Gbps & 54{,}000 & 13{,}500 & N/A \\
\bottomrule
\multicolumn{8}{l}{\footnotesize $^\dagger$S2 reports the final tracking-mode
  equilibrium under continuous warm-start EGT stepping from Phase~I.}\\
\multicolumn{8}{l}{\footnotesize \phantom{$^\dagger$}Transient Iters denotes
  the per-step tracking iteration count, not a single cold-start convergence.}\\
\end{tabular}
\end{table*}

\begin{figure*}[t]
    \centering
    \includegraphics[width=0.95\textwidth]{multi_scenario_1000ue}
    \caption{Emulation results for 1,000 UEs across four scenarios:
    (a) population fractions converging to Nash equilibrium,
    (b) end-to-end delays versus PDB thresholds, and
    (c) throughput and QoS violations.}
    \label{fig:results}
\end{figure*}

%--------------------------------------------------------------
\subsection{Key Analysis and Findings}
\label{sec:key_findings}
%--------------------------------------------------------------

\textbf{Nash equilibrium and delay equalisation.}
In S1, the EGT controller converges to
$x_\text{mec}^\text{eMBB} = 0.156$ and
$x_\text{mec}^\text{URLLC} = 0.222$, routing 12.9~Gbps to the
edge node and 58.6~Gbps to the central cloud. At this equilibrium,
the eMBB end-to-end delays are 27.9~ms (MEC) and 17.9~ms (CC),
yielding a weighted average of 19.5~ms. URLLC delays are 19.5~ms
(MEC) and 17.9~ms (CC), weighted average 18.3~ms. The residual
$\Delta d = 10.0$~ms gap between MEC and CC at the stopping
criterion is the expected numerical consequence of the payoff
convergence tolerance $\varepsilon = 0.01$, not a failure of
the replicator dynamics.

\textbf{eMBB QoS compliance.}
Zero eMBB violations are recorded in S1, S3, and S4, confirming
that the 19.5~ms weighted average delay is well within the
300~ms eMBB PDB. The EGT controller effectively steers eMBB
traffic toward $\text{UPF}_\text{cc}$ during load surges,
preserving the edge node for delay-sensitive traffic. The 10:1
cloud-to-edge processing exponent ratio ($k_\text{mec}/k_\text{cc}
= 10$) drives sufficient eMBB spillover to the central cloud to
prevent the edge node from saturating for this slice class.

\textbf{URLLC infrastructure ceiling.}
At the S1 peak equilibrium, both MEC and CC delays
(19.5~ms and 17.9~ms respectively) exceed the
10~ms URLLC PDB. This results in 26,316 cumulative
URLLC UE-step violations over Phase~II.
This is not a policy failure: the EGT controller achieves
the optimal load distribution given the hardware, but the
total offered URLLC load at $5.5\times$ multiplier
exceeds what the deployed nodes can serve within 10~ms.
Specifically, evaluating Eq.~(3a) for the URLLC group at 10~ms reveals an absolute edge capacity ceiling:
\begin{equation*}
  k_\text{mec} \cdot (\text{Load}) > \ln(10 - 0.25) \implies \text{Load} \approx 4.7 \text{ Gbps}
\end{equation*}
Because the total URLLC load is 27.5~Gbps, the hard infrastructure ceiling of 4.7~Gbps per slice cannot be overcome by steering algorithms alone.
Future capacity planning must account for this boundary explicitly.

\textbf{Transient recovery and numerical $\varepsilon$-boundaries (S3, S4).}
Although S1, S3, and S4 share identical load profiles and theoretically converge to the same unique interior Nash equilibrium, the discrete controller terminates at the $\varepsilon=0.01$ stopping boundary. Because S1 approaches this boundary from $x=0.5$ while S4 approaches it from $x=0.05$, they settle on slightly different boundary edge fractions (0.156 vs 0.128 respectively).

The transient behaviour is highly dependent on initialisation. S4 (CC-biased initialisation, 5\% MEC) converges rapidly in 270~iterations but accumulates a staggering 44,027 URLLC violations. This occurs because its starting state (95\% load on the central cloud) drives initial CC delays past 27~ms, ensuring virtually all URLLC sessions violate the 10~ms PDB until the traffic can be migrated. S3 requires 670~iterations
to recover. The longer S3 recovery reflects
the exponential delay penalty incurred when the edge node is
overloaded from a sub-optimal starting policy, confirming that
EGT is globally convergent but that convergence speed depends
on initial conditions.

\textbf{S2: extreme spike, numerical duality, and dynamic tracking.}
Scenario S2 ($10\times$ load, 130~Gbps total) reveals an important
theoretical property of the EGT replicator dynamics under sudden
extreme load transitions: a \emph{mathematical duality} between
cold-start single-shot optimisation and continuous dynamic tracking.

When EGT is initialised from a neutral state ($x = 0.5$) and
converged in a single call at $10\times$ load with the standard
tolerance $\varepsilon = 0.01$, the optimiser locks into a
sub-optimal numerical trap at $x_\text{mec}^\text{eMBB} \approx 0.380$.
This overloads the edge node, producing catastrophic MEC processing
delay ($d_\text{eMBB}^\text{MEC} \approx 2.4\times10^6$~ms) and a
MEC load of 49.4~Gbps. The true Nash equilibrium at $10\times$
load (reachable only by tightening $\varepsilon$ to $10^{-5}$,
requiring 7,673 iterations) lies at
$x_\text{mec}^\text{eMBB} \approx 0.136$ with an average delay of 188.8~ms.

In the actual stadium simulation, however, EGT operates in
\emph{continuous warm-start tracking} mode: it steps through
the load trajectory from Phase~I, so that when Phase~II begins,
the population state is already close to the active equilibrium
($x_\text{mec}^\text{eMBB} \approx 0.229$ at Phase~I end). EGT
progressively migrates traffic across the 135 Phase~II timesteps,
stabilising at $x_\text{mec}^\text{eMBB} = 0.147$ with an average
eMBB delay of 189.3~ms and a MEC load of 23.8~Gbps. Cumulative
eMBB violations total 13,491, representing only the transient
steps during active migration; once EGT converges, MEC and CC
delays fall below 300~ms and violations cease. This demonstrates
that \emph{continuous dynamic step-by-step EGT tracking with warm-start prevents the local numerical traps common in single-shot
static optimisation under sudden load transitions}, and constitutes
a practically important empirical observation of this work.

\textbf{Baseline comparison.}
Under static 50/50 routing, 35.8~Gbps is steered to the MEC node,
driving average eMBB MEC delay to $\approx 4.15\times10^4$~ms and
average URLLC MEC delay to $769$~ms. The central cloud, receiving
an equal 35.8~Gbps, remains lightly loaded at 6.6~ms. The
weighted average eMBB delay across both UPFs is 20,736~ms and
the URLLC average is 388~ms, both far exceeding their respective
PDB targets. This confirms that static routing catastrophically
overloads the edge node while leaving the cloud idle, whereas
the EGT controller stabilises both nodes at a jointly optimal
operating point.

\section{Future Work}
\label{sec:future}

The validated EGT steering logic opens several avenues
for future research.

\subsection{Predictive ML-Based Reconfiguration Triggers}

A critical enhancement is the transition from periodic
scheduling to an event-driven reconfiguration trigger.
We propose the integration of a Bidirectional LSTM
(BLSTM) model within the EGT control loop. Trained on
mobility-driven QoS degradation traces, this model
would utilise three primary indicators---the
reconfiguration countdown $T_g^{\text{reconf}}$, the
predicted session violation time $T_{g}^{\text{vio}}$,
and the instantaneous violation count $L_a$---to
proactively initiate EGT re-convergence before
violations occur. This is expected to reduce
reconfiguration overhead and maintain higher QoS
compliance, following the approach of~\cite{LeyvaUPC2023}.

\subsection{Economic and OPEX-Aware Optimisation}

Future work will address the operational cost dimension
of edge steering. MEC resources are typically more
expensive to maintain than centralised cloud infrastructure.
We plan to extend the payoff function $u(\mathbf{x})$
to include an economic weight factor that penalises
edge strategies for non-critical traffic, enabling
operators to reserve premium MEC capacity for URLLC
slices while directing eMBB traffic to more
cost-effective cloud resources during low-congestion
periods.

\subsection{Sustainability and Energy-Aware Steering}

In alignment with Green-6G objectives, we intend to
incorporate power consumption models into the steering
logic. Since distributed edge UPFs can consume more
energy per bit than consolidated cloud data centres, an
energy-aware orchestration layer could adjust the EGT
equilibrium to enable reduced-power states for local
UPFs during the post-event dispersal phase, balancing
latency performance against the venue's overall energy
footprint.

\subsection{Validation on Physical Testbeds}

Further experimental validation is required using
NSSF-enabled multi-UPF deployments in physical testbeds,
with real over-the-air radio rather than emulated UE
sessions. Investigating the impact of control-plane
signalling overhead during mass PDU session events
remains a high priority for hardware-in-the-loop
studies.

\section{Conclusion}

This paper has presented a slice-aware, multi-UPF
selection framework for 5G stadium environments,
integrating Evolutionary Game Theory (EGT) replicator
dynamics with an OpenAirInterface (OAI) emulation
testbed supporting 1,000 concurrent PDU sessions. A
high-performance Go-based user-plane multiplexing engine
enabled 100\% PDU session establishment success at scale,
confirming the feasibility of high-density emulation on
standard hardware.

The key findings are threefold. First, the EGT-based
steering achieves zero eMBB QoS violations across all
standard and overload scenarios by exploiting the
10:1 cloud-to-edge processing exponent ratio, confirming
the practical effectiveness of population-level steering
at scale. Second, the EGT controller converges to a stable 
equilibrium across all evaluated scenarios. While static 
cold-starts under extreme load shifts require up to 670 
iterations, continuous warm-start tracking resolves in 1 
iteration under standard load transitions, and up to 147 iterations 
per step during extreme $10\times$ load surges. Third, and importantly,
the study reveals a hard infrastructure boundary for
URLLC QoS compliance: during the 5.5$\times$ halftime
surge, edge-node processing capacity is saturated
regardless of the steering policy applied. Using the Phase II 
window of 27,000 URLLC UE-steps as the denominator, the 26,316 
URLLC violations in S1 represent a catastrophic 97.5\% violation rate. 
This demonstrates that under severe congestion, slice-agnostic delay 
minimisation cannot protect strict PDB slices from infrastructure saturation.

Future research will explore slice-margin payoff adjustments, BLSTM-based
predictive triggers to reduce reactive reconfiguration
overhead, economic cost weighting in the payoff function,
and energy-aware steering for post-event dispersal phases.

\section*{Acknowledgment}
The authors would like to thank the OpenAirInterface Software
Alliance for providing the core networking framework used in
this research.

\bibliographystyle{IEEEtran}
\bibliography{references}

\end{document}
